\documentclass[slidestop,compress,mathserif]{beamer}
%\documentclass[slidestop,compress,mathserif,handout]{beamer}

%\documentclass[xcolor=dvipsnames,handout]{beamer}
%\documentclass[xcolor=dvipsnames]{beamer}

%\documentclass[handout]{beamer}

%%% To get rid of solutions on handouts:
\newcommand{\soln}[1]{\textit{\textcolor{darkGray}{#1}}}				% For slides
%\newcommand{\soln}[1]{ }	% For handouts

% to get pausing to work properly on slides
\newcommand{\hide}[1]{#1}	% For slides
%\newcommand{\hide}[1]{ }	% For handouts
\usepackage{tikz}


%\usepackage{multicol}
\usepackage{amsfonts}
%\usepackage[pdftex,dvipsnames]{color}
\usepackage{graphicx}
\usepackage{subfigure}
%\usepackage{picinpar}
\usepackage{pifont}
\usepackage{pgf,pgfarrows,pgfnodes}
%\usepackage{wasysym,manfnt,phaistos,empheq}
\usepackage[english]{babel}
\usepackage{pgfpages}
\usepackage{natbib}
\usepackage{hyperref}
\usepackage{multimedia}
%\usepackage{amsfonts,amstext,amssymb,amsbsy,amsopn,amsthm,eucal,latexsym,mathrsfs}
\usepackage{amsmath,amsfonts,amstext,amssymb,amsbsy,amsopn,amsthm,eucal,latexsym,mathrsfs}
\usepackage{ulem}
\usepackage{setspace}
\usepackage{array}
%\usepackage{rotating}
\usepackage{multirow}
\usepackage{verbatim}
\usepackage{multicol}

\setbeamertemplate{navigation symbols}{}

%\usepackage{tikz}
%\usetikzlibrary{arrows,shapes,trees,backgrounds}


%\setbeameroption{show notes on second screen}
%\setbeameroption{show notes}
%\setbeameroption{show only notes}

\definecolor{links}{HTML}{2A1B81}
\hypersetup{colorlinks,linkcolor=,urlcolor=links}

\newtheorem*{principle}{Inscrutibility Principle}
\newtheorem*{punchline}{Punch Line}
\newtheorem{defn}{Definition}

\definecolor{Scarlet}{RGB}{140,17,17}
\definecolor{VassarRed}{RGB}{128,0,0}

% "dinglist" environment
  \renewenvironment{dinglist}[2][blue]
  {\begin{list}{\textcolor{blue}{\ding{#2}}}{}}{\end{list}}
  % Symbol definitions for these lists
  \newcommand{\DingListSymbolA}{43}
  \newcommand{\DingListSymbolB}{243}
  \newcommand{\DingListSymbolC}{224}
  \newcommand{\DingListSymbolD}{219}
  \newcommand{\DingListSymbolCheck}{52}
  \newcommand{\DingListSymbolCross}{56}


  \newenvironment{ballotenv}
{\only{%
\setbeamertemplate{itemize item}{\ding{45}}%
\setbeamertemplate{itemize subitem}{\ding{46}}%
\setbeamertemplate{itemize subsubitem}{\ding{46}}}} {}
\setbeamertemplate{itemize item}{\ding{49}}
\setbeamertemplate{itemize subitem}{\ding{47}}
\setbeamertemplate{itemize subsubitem}{\ding{47}}


%User defined colors: See colors section
\xdefinecolor{oiBlue}{rgb}{0.15, 0.35, 0.55}
\xdefinecolor{gray}{rgb}{0.5, 0.5, 0.5}
\xdefinecolor{darkGray}{rgb}{0.3, 0.3, 0.3}
\xdefinecolor{darkerGray}{rgb}{0.2, 0.2, 0.2}
\xdefinecolor{rubineRed}{rgb}{0.89,0,0.30}
\xdefinecolor{linkCol}{rgb}{0.11,0.49,0.95}	
\xdefinecolor{irishGreen}{rgb}{0,0.60,0}	
\xdefinecolor{darkturquoise}{rgb}{0.44, 0.58, 0.86}
\definecolor{lightGreen}{rgb}{0.533,0.765,0.42}
\xdefinecolor{Regalia}{HTML}{522D80}
\xdefinecolor{ClemsonOrange}{HTML}{EA6A20}

\definecolor{duke@LightGrey}{RGB}{200,200,200}\definecolor{DarkGreen}{RGB}{0,100,0}
\definecolor{Oranges}{RGB}{255,127,0}
\definecolor{LightGray}{RGB}{211,211,211}

%\setbeamertemplate{footline}{%
%  \raisebox{5pt}{\makebox[\paperwidth]{\hfill\makebox[10pt]{\scriptsize\insertframenumber}}}}

\setbeamercolor{equation background}{fg=black,bg=duke@LightGrey}
  % Boxed equation
  \newcommand{\eqbox}[2][0.6]{%
  \centerline{
  \begin{beamerboxesrounded}[lower=equation background,width=#1\hsize,shadow=true]{}
\parbox{#1\hsize}{%
      \[
        \textcolor{black} {#2}
      \]}
  \end{beamerboxesrounded}
}}

\AtBeginSection[] {
  \begin{frame}<beamer>\frametitle{Outline}
    \tableofcontents[currentsection,hideothersubsections]
  \end{frame}
}
%
%
%\AtBeginSubsection[] {
%  \begin{frame}<beamer>\frametitle{Outline}
%    \tableofcontents%[currentsection,currentsubsection]
%  \end{frame}
%}

%\usecolortheme[RGB={82,45,128}]{structure}
%\usecolortheme[RGB={162,80,22}]{structure}
\usecolortheme[RGB={128,0,0}]{structure}
\usetheme[secheader]{Boadilla}
%\usetheme[height=7mm]{Rochester}
%\usetheme{Copenhagen}
%\usetheme{Antibes}
%\usecolortheme{seahorse}
%\usecolortheme{crane}
%\usecolortheme{rose}
%\usefonttheme[onlylarge]{structurebold}
%\usefonttheme[onlymath]{serif}



\def\diag{{\rm diag}}


\def\E{\mathbb{E}}
\def\Prob{\mathbb{P}}
\def\argmin{{\rm argmin}}
\def\argmax{{\rm argmax}}
\def\Def{\stackrel{def}{=}}


\newtheorem{assumption}{Assumptions}
\newtheorem*{proposition}{Proposition}
\newtheorem*{remark}{Remark}



%\setbeamercolor{disc title}{bg=oiBlue!40!white!60,fg=blue}
\setbeamercolor{disc body}{bg= Regalia!20!white!80,fg= Regalia!80!black!90}

\setbeamercolor{clicker ungraded title}{bg=irishGreen!80!white!50,fg=irishGreen!30!black!90}
\setbeamercolor{clicker ungraded body}{bg=irishGreen!20!white!80,fg=irishGreen!30!black!90}

\setbeamercolor{clicker review title}{bg=gray!80!white!80,fg=oiBlue!80!black!90}
\setbeamercolor{clicker review body}{bg=gray!30!white!90,fg=oiBlue!80!black!90}

\setbeamercolor{code body}{bg=gray!20!white!80,fg=black}


% Custom commands
\newcommand{\degree}{\ensuremath{^\circ}}
\newcommand{\Note}[1]{
\rule{2.5cm}{0.25pt} \\ \textit{\scriptsize {\textcolor{rubineRed}{Note:} \textcolor{gray}{#1}}}}
\newcommand{\ct}[1]{
\vfill
{\tiny #1}}
\newcommand{\Remember}[1]{\textit{\scriptsize{\textcolor{orange}{Remember:} \textcolor{gray}{#1}}}}
\newcommand{\red}[1]{\textit{\textcolor{rubineRed}{#1}}}
\newcommand{\pink}[1]{\textit{\textcolor{rubineRed!90!white!50}{#1}}}
\newcommand{\green}[1]{\textit{\textcolor{irishGreen}{#1}}}
\newcommand{\webURL}[1]{\urlstyle{same}\textit{\textcolor{linkCol}{\url{#1}}} }
\newcommand{\webLink}[2]{\href{#1}{\textcolor{linkCol}{{#2}}}}
\newcommand{\mail}[1]{\href{mailto:#1}{\textit{\textcolor{linkCol}{#1}}}}
\newcommand{\hl}[1]{\textit{\textcolor{oiBlue}{#1}}}
\newcommand{\hlGr}[1]{\textit{\textcolor{lightGreen}{#1}}}
\newcommand{\mathhl}[1]{\textcolor{oiBlue}{\ensuremath{#1}}}
\newcommand{\ex}[1]{\textcolor{blue}{{{\small (#1)}}}}
\newcommand{\disc}[1]{
\begin{beamerboxesrounded}[shadow = true, lower = disc body, upper = disc title]{}
#1
\end{beamerboxesrounded}
}

\newcommand{\cl}[1]{
\begin{beamerboxesrounded}[shadow = true, lower = clicker ungraded body, upper = clicker ungraded title]{Question}
$\:$ \\
#1
\end{beamerboxesrounded}
}

\newcommand{\clR}[1]{
\begin{beamerboxesrounded}[shadow = true, lower = clicker review body, upper = clicker review title]{\red{Review question} }
$\:$ \\
#1
\end{beamerboxesrounded}
}

\newcommand{\formula}[2]{
\begin{beamerboxesrounded}[shadow = true, lower = white, upper = clicker review body]{#1}
#2
\end{beamerboxesrounded}
$\:$ \\
}

\newenvironment{twocol}[4]{
\begin{columns}[c]
\column{#1\textwidth}
#3
\column{#2\textwidth}
#4
\end{columns}
}


\newenvironment{slot}[2]{
\begin{array}{c}
\underline{#1} \\
#2
\end{array}
}

\newcommand{\pr}[1]{
\left( #1 \right)
}

\newcommand{\solnMult}[1]{
\item[] \vspace{-0.59cm}
\only<beamer| beamer:1>{\item #1}
\soln{\only<2->{\item \red{#1}}}
}

%\newcommand{\codechunk}[1]{
%\begin{beamerboxesrounded}[shadow = true, lower = code body]{}
%{\small #1}
%\end{beamerboxesrounded}
%}

% Change margin

\newenvironment{changemargin}[2]{%
\begin{list}{}{%
\setlength{\topsep}{0pt}%
\setlength{\leftmargin}{#1}%
\setlength{\rightmargin}{#2}%
\setlength{\listparindent}{\parindent}%
\setlength{\itemindent}{\parindent}%
\setlength{\parsep}{\parskip}%
}%
\item[]}{\end{list}}

% Footnote

\long\def\symbolfootnote[#1]#2{\begingroup%
\def\thefootnote{\fnsymbol{footnote}}\footnote[#1]{#2}\endgroup}

% Commands from the book
\newenvironment{data}[1]{\texttt{#1}}{}
\newenvironment{var}[1]{\texttt{#1}}{}
\newenvironment{resp}[1]{\texttt{#1}}{}






%%%%%%%%%%%%%%%%%%%%%%%%%%%%%%%%%%%%%%%%%%%%%%%%%%%%%%%%%%%%%%%%%%%%%%%%%%%%%%%%%%%%%%%%%%%%%%%

\title[Chapter 8]{Chapter 8}
\subtitle{Limit Theorems}

%%%%%%%%%%%%%%%%%%%%%%%%%%%%%%%%%%%%%%%%%%%%%%%%%%%%%%%%%%%%%%%%%%%%%%%%%%%%%%%%%%%%%%%%%%%%%%%


\author[Jingchen (Monika) Hu] % (optional, use only with lots of authors)
{Jingchen (Monika) Hu}
% - Give the names in the same order as the appear in the paper.
% - Use the \inst{?} command only if the authors have different
%   affiliation.

\institute[Vassar] % (optional, but mostly needed)
{Vassar College}
% - Use the \inst command only if there are several affiliations.
% - Keep it simple, no one is interested in your street address.

\date[MATH 241] % (optional, should be abbreviation of conference name)
{MATH 241}
% - Either use conference name or its abbreviation.
% - Not really informative to the audience, more for people (including
%   yourself) who are reading the slides online

\subject{MATH 241}
% This is only inserted into the PDF information catalog. Can be left
% out.



% If you wish to uncover everything in a step-wise fashion, uncomment
% the following command:

%\beamerdefaultoverlayspecification{<+->}


\begin{document}

%\begin{frame}%[plain]
%\includegraphics[width = \textwidth]{figures/2015DukeNCAA}
%\end{frame}

%{ % all template changes are local to this group.
%\addtocounter{framenumber}{-1}
%    \setbeamertemplate{navigation symbols}{}
%    \begin{frame}[plain]
%        \begin{tikzpicture}[remember picture,overlay]
%            \node[at=(current page.center)] {
%                \includegraphics[width=1.25\paperwidth]{figures/2015DukeNCAA}
%            };
%        \end{tikzpicture}
%     \end{frame}
%}

%%%%%%%%%%%%%%%%%%%%%

% Title Page

\begin{frame}%[plain]
\titlepage
\end{frame}


%%%%%%%%%%%%%%%%%%%%%%
%\addtocounter{framenumber}{-1}
%
%\begin{frame}\frametitle{Annoucement}
%
%\begin{itemize}
%%\item HW7: \red{due now!}
%\item HW8: \red{due Tuesday, Nov 20th}
%
%
%\vspace{0.5cm}
%\item Course evaluation open tomorrow.
%  \begin{itemize}
%  \item Activate between 11/14 - 12/3.
%  \item If response rate $\geq 80\%$, drop the lowest quiz.
%  \end{itemize}
%
%\vspace{0.5cm}
%\item Next quiz: Tuesday, Nov 18th\\
%  \begin{itemize}
%  \item Topic: function of a continuous RV, i.e., find $f_Y(y)$ where $Y = g(X)$.
%  %\item To prepare: do homework questions in Chapter 5: 37, 39, 40, TE29
%  \end{itemize}
%
%
%%\item Midterm: Tuesday, Feb 25th
%%\begin{itemize}
%%\item Close book, in class exam (75 min)
%%\item ONE page cheat sheet {\bf made by yourself} (A4 size)
%%\item Calculators are allowed, but not cell phones, tablets or laptops
%%\end{itemize}
%
%\end{itemize}
%
%
%\end{frame}
%

%
%%%%%%%%%%%%%%%%%%%%%%
%\begin{frame}{Outline}
%%\tableofcontents[hideallsubsections,pausections]
%\tableofcontents[hideallsubsections]
%\end{frame}

%%%%%%%%%%%%%%%%%%%%%%%%%%%%%%%%%%%%%%%%%%%
%\begin{frame}\frametitle{Recap}
%
%
%\end{frame}





%%%%%%%%%%%%%%%%%%%%%
%\begin{frame}{Outline}
%%\tableofcontents[hideallsubsections,pausections]
%\tableofcontents[hideallsubsections]
%\end{frame}

%%%%%%%%%%%%%%%%%%%%%%%%%%%%%%%%%%%%%%%%%%%
%\section{Introduction}
%%%%%%%%%%%%%%%%%%%%%%%%%%%%%%%%%%%%%%%%%%%
%\begin{frame}\frametitle{Overview}
%
%\begin{itemize}
%\item The most important limit theorems are classified under two headings: ``laws of large numbers" and ``central limit theorems"
%\vspace{5mm}
%\pause
%\item Theorems under ``laws of large numbers" (LLN)
%\begin{itemize}
%\item Stating conditions under which the average of a sequence of random variables converges (in some sense) to the expected average
%\item We will study the strong LLN
%\end{itemize}
%\pause
%\vspace{5mm}
%\item Theorems under ``central limit theorems" (CLT)
%\begin{itemize}
%\item Determining conditions under which the sum of a large number of random variables has a probability distribution that is approximately normal
%\end{itemize}
%\end{itemize}
%\end{frame}
%
%%%%%%%%%%%%%%%%%%%%%%%%%%%%%%%%%%%%%%%%%%%
%\section{Markov's and Chebyshev's inequalities}
%%%%%%%%%%%%%%%%%%%%%%%%%%%%%%%%%%%%%%%%%%%
%\begin{frame}%\frametitle{Chebyshev's inequality}
%
%\begin{block}{Chebyshev's inequality}
%If $X$ is a random variable with finite mean $\mu$ and variance $\sigma^2$, then for any value $k > 0$
%$$P\{|X - \mu| \geq k\} \leq \frac{\sigma^2}{k^2}$$
%\end{block}
%\pause
%{\small{
%To prove Chebyshev's inequality, useful proposition: Markov's inequality\\
%If $X$ is a random variable that takes only nonnegative values, then for any value $a > 0$
%$$P\{X \geq a\} \leq \frac{E[X]}{a}$$
%}}
%
%
%\end{frame}
%
%\begin{frame}\frametitle{Chebyshev's inequality applied to known distributions}
%\begin{itemize}
%\item Earlier we saw we can use Chebyshev's inequality and/or Markov' inequality to derive bounds on probabilities
%\item Chebyshev's inequality: $P\{|X - \mu| \geq k\} \leq \frac{\sigma^2}{k^2}$
%\item How about for known distributions?
%\pause
%\end{itemize}
%{\small{
%\disc{If $X$ is uniformly distributed over the interval $(0, 10)$. What is the exact probability of $P\{|X - 5| > 4\}$ and the upper bound of this probability obtained by Chebyshev's inequality?}
%\pause
%\begin{itemize}
%\item Exact probability:
%$$P\{|X - 5| > 4\} = .20$$
%\pause
%\item Upper bound by Chebyshev's inequality ($E[X] = 5, Var(X) = \frac{25}{3}$):
%$$P\{|X - 5| > 4\} \leq \frac{25}{3(16)} \approx .52$$
%\pause
%\item Chebyshev's inequality gives very rough/not close (but correct!) bound
%\end{itemize}
%}}
%\end{frame}
%
%
%\begin{frame}\frametitle{Chebyshev's inequality applied to known distributions}
%\cl{If $X$ is a normal random variable with mean $\mu$ and variance $\sigma^2$. What is the exact probability of $P\{|X - \mu| > 2\sigma\}$ and the upper bound of this probability obtained by Chebyshev's inequality?}
%\pause
%{\small{
%\begin{itemize}
%\item Exact probability:
%$$P\{|X - \mu| > 2\sigma\} = P\{|\frac{X - \mu}{\sigma}| > 2\} = 2[1 - \Phi(2)] \approx .0456$$
%\pause
%\item Upper bound by Chebyshev's inequality:
%$$P\{|X - \mu| > 2\sigma\} \leq \frac{1}{4}$$
%\pause
%\item Again, Chebyshev's inequality gives very rough/not close (but correct!) bound
%\end{itemize}
%}}
%\end{frame}
%
%\begin{frame}\frametitle{Usefulness of Chebyshev's inequality}
%\begin{itemize}
%\item Chebyshev's inequality is valid for all distributions of the random variable $X$
%\vspace{5mm}
%\item Do not expect the bound on the probability to be very close to the actual probability in most cases
%\vspace{5mm}
%\pause
%\item However, if we do not know the exact probability distribution but its mean and variance are known, Chebyshev's and Markov's inequalities are very useful to derive bounds on probabilities
%\vspace{5mm}
%\pause
%\item On the other hand, if the exact distribution is known, we won't need the inequalities to bound the probabilities
%\end{itemize}
%\end{frame}
%
%\begin{frame}
%\disc{From past experience, a professor knows that the test score of a student taking her final exam is a random variable with mean $75$. (a) Give an upper bound for the probability that a student's test score will exceed $85$, suppose that the variance is $25$. (b) What can be said about the probability that a student will score between $65$ and $85$?}
%\pause
%{\small{
%\begin{itemize}
%\item[(a)] Using Markov's inequality
%$$P\{X \geq 85\} \leq \frac{E[X]}{85} = \frac{75}{85} = \frac{15}{17}$$
%\pause
%\item[(b)] Using Chebyshev's inequality
%$$P\{65 \leq X \leq 85\} = P\{|X - 75| \leq 10\} = 1 - P\{|X - 75| \geq 10\} \geq 1 - \frac{25}{100} = \frac{3}{4}$$
%\end{itemize}
%}}
%\end{frame}
%
%%\begin{frame}
%%\begin{block}{The weak law of large numbers}
%%Let $X_1, X_2, ..., X_n$ be a sequence of independent and identically distributed random variables, each having finite mean $E[X_i] = \mu$. Then, for any $\epsilon > 0$, as $n \to \infty$
%%$$P\{|\frac{X_1 + ... + X_n}{n} - \mu| \geq \epsilon \} \to 0$$
%%\end{block}
%%\pause
%%Proof (assuming that the random variables have a finite variance $\sigma^2$)
%%\end{frame}

%%%%%%%%%%%%%%%%%%%%%%%%%%%%%%%%%%%%%%%%%%
\section{Central limit theorem}
%%%%%%%%%%%%%%%%%%%%%%%%%%%%%%%%%%%%%%%%%%
\begin{frame}
\begin{block}{The central limit theorem}
Let $X_1, X_2, ...$ be a sequence of independent and identically distributed random variables, each having mean $\mu$ and variance $\sigma^2$. Then the distribution of
$$\frac{X_1 + X_2 + ... + X_n - n\mu}{\sigma \sqrt{n}}$$
tends to the standard normal as $n \to \infty$. 
%That is for $-\infty < a < \infty$, as $n \to \infty$
%$$P\{\frac{X_1 + X_2 + ... + X_n - n\mu}{\sigma \sqrt{n}} \leq a\} \to \frac{1}{\sqrt{2\pi}}\int_{-\infty}^{a}e^{-\frac{x^2}{2}}dx$$
\end{block}

\vspace{5mm}
In words, if the random variables have a finite mean $\mu$ and a finite variance $\sigma^2$, then the distribution of the sum of the first $n$ of them is, for large $n$, approximately that of a normal random variable with mean $n\mu$ and variance $n\sigma^2$.
\end{frame}



\begin{frame}\frametitle{Applications of the CLT}
\cl{A person has $100$ light bulbs whose lifetimes are independent exponentials with mean $5$ hours. If the bulbs are used one at a time, with a failed bulb being replaced immediately by a new one. Suppose that it takes a random time, uniformly distributed over $(0, .5)$, to replace a failed bulb. Approximate the probability that all bulbs have failed by time 550. Note that if $X \sim \textrm{Exponential}(\lambda)$, then $E[X] = \frac{1}{\lambda}$ and $Var(X) = \frac{1}{\lambda^2}$.}

\end{frame}

\begin{frame}

{\small{
\begin{itemize}
\item Let $X_i$ be the lifetime of bulb $i$ and $R_i$ be the time to replace bulb $i$
\pause
\item The desired probability is $P\{\sum_{i=1}^{100}X_i + \sum_{i=1}^{99}R_i \leq 550\}$
\item What is the distribution of $\sum_{i=1}^{100}X_i + \sum_{i=1}^{99}R_i$?
\pause
\item Sum of two independent normal is normal with $\mu_1+\mu_2$ and $\sigma_1^2 + \sigma_2^2$
\pause
\item Thus $\sum_{i=1}^{100}X_i + \sum_{i=1}^{99}R_i$ has mean $100\times 5 + 99\times .25 = 524.75$ and variance $2500 + 99/48 = 2502$
\item Then the desired probability is approximately equal to
$$P\{N(0, 1) \leq \frac{550 - 524.75}{\sqrt{2502}}\} = P\{N(0, 1) \leq .505\} = .693$$
\end{itemize}

}}
\end{frame}
%
%
%%%%%%%%%%%%%%%%%%%%%%%%%%%%%%%%%%%%%%%%%%%
%\section{The strong law of large numbers}
%%%%%%%%%%%%%%%%%%%%%%%%%%%%%%%%%%%%%%%%%%%
%\begin{frame}
%\begin{block}{The strong law of large numbers}
%Let $X_1, X_2, ...$ be a sequence of independent and identically distributed random variables, each having a finite mean $\mu = E[X_i]$. Then, with probability $1$, as $n \to \infty$
%$$\frac{X_1 + X_2 + ... + X_n}{n} \to \mu$$
%\end{block}
%\pause
%\begin{itemize}
%\item The strong LLN states that $P\{\underset{n\to \infty}{\lim}\frac{X_1 + ... + X_n}{n} = \mu\} = 1$
%\item In words, the average of a sequence of independent random variables having a common distribution will with probability $1$, converge to the mean of that distribution
%\end{itemize}
%\end{frame}
%
%
%\begin{frame}\frametitle{The strong LLN and the interpretation of probability}
%Back in Chapter 2, we discussed two different schools of interpreting probability
%\begin{itemize}
%\item \red{Frequentist} interpretation: The probability of event $A$ is the proportion of times (frequency) that A occurs in an infinite sequence (or very long run) of separate tries of the experiment.
%\[ P(A)= \lim_{n \rightarrow \infty} \frac{\mbox{\# times A happens}}{n}.\]
%\pause
%\item A {\red{Bayesian}} can pick whatever number they prefer for $P(A)$, based on their own personal experience and intuition, provided that number is consistent with all of the other probabilities they choose in life.
%\pause
%\item The strong LLN can help understand the frequentist interpretation
%\end{itemize}
%\end{frame}
%
%
%\begin{frame}{Application of the strong LLN}
%\begin{itemize}
%\item Suppose that a sequence of independent trials of some experiment is performed
%\item Let $A$ be a fixed event of the experiment, and denote by $P(A)$ the probability that $A$ occurs on any particular trial
%\pause
%\item Letting
%\[ X_i = \begin{cases}
%    1 & \text{if $A$ occurs on the $i$-th trial} \\
%    0& \text{if $A$ does not occur on the $i$-th trial}
%\end{cases}
% \]
%\pause
%By the strong LLN, with probability $1$
%$$\frac{X_1 + ... + X_n}{n} \to E[X] = P(A)$$
%\pause
%\item $X_1 + ... + X_n$ represents the number of times that the event $A$ occurs in the first $n$ trials, the limiting proportion of time that the event $A$ occurs is just $P(A)$ with probability $1$
%\end{itemize}
%
%\end{frame}
\end{document}
